\title{Inflation Dynamics and Relative Price Dispersion in South Africa: Implications for Optimal Inflation Targeting}
\author {Keagile Lesame\footnote{Corresponding author. National Treasury, South Africa. Email: Keagile.Lesame@treasury.gov.za} \and
Xolani Sibande\footnote{Economic Research Department, South African Reserve Bank. Email: xolani.sibande@resbank.co.za.}}


\date{\today}
\maketitle

\begin{abstract}
This paper investigates the dynamic relationship between inflation and relative price dispersion (RPD) in South Africa, focusing on the period from 2009 to 2025, a timeframe marked by significant structural changes in inflation targeting. Drawing on both theoretical models—the menu cost and signal extraction frameworks—the study tests their predictions on how inflation influences RPD using a semi-parametric approach. It extends the analysis by employing quantile-on-quantile regression to capture state-dependent effects across the distribution of relative price variability. The South African case is particularly insightful due to its evolving inflation targets, transitioning from a 3-6 percent band (2000) to a 4.5 percent midpoint target (2017), and moving towards a 3 percent target projected for 2025. These policy shifts offer a unique opportunity to explore the nonlinear, U-shaped nature of the inflation-RPD relationship, challenging the traditional view that zero inflation is optimal. Empirically, this study contributes to the scant emerging market literature, complementing prior work but incorporating inflation target changes. The findings have important implications for monetary policy design in emerging economies aiming to lower inflation targets while minimizing distortive effects on price signals and resource allocation.
\end{abstract}

\noindent\textbf{Keywords}: Relative Price Dispersion, Inflation, Monetary Policy\\
\textbf{JEL Codes}: E31, E52, E58
\newpage